% Technical Background, section 1 (written 2026-11-06; replaces the parked logistic/ridge/1D-CNN
%   plan, now retired/05_risk_models_parked.tex).
% Pointed back to by: methodology/02_models (NOT WRITTEN), the base model of the who analysis.
% Model: Aida 3.2 / 3.2.1, the method in the section body and one subsection one step deeper,
%   opening on the parent's limitation.
% Does NOT decide the open base choice (recalibrated SCORE2 versus refitted SCORE2 factors, root
%   CLAUDE.md "Open"): both are presented as general methods; the choice is Methodology.
% Hazard closer (guide 3.6): an added block can look informative by correcting the weights of a
%   published score. Own reasoning, not cited; see Hlatky 2009 (AHA) in
%   ~/Desktop/Speciale/CGPS_statistics-/reports/Geometry added value over SCORE2.md if a source is wanted.
% Notation: z_i, gamma_0, gamma_k, LP_i, g, a, gamma. Avoids beta, eta, b, S, j (the GLMM) and the
%   FPCA symbols.
% LEFT OUT on purpose (Methodology): which link, which predictors, SCORE2-OP, follow-up time,
%   degrees of freedom and the Riley budget.

\section{Clinical Risk Prediction Models}
\label{sec:tech-risk-models}

A clinical risk prediction model estimates the absolute probability that a person with given characteristics develops an outcome within a stated time horizon \cite{steyerberg2019cpm}.
Let $\mathbf{z}_i = (z_{i1}, \dots, z_{iP})^\top$ denote the predictors of person $i$.
The model combines them into a linear predictor and maps it to a risk through a link function $g$,
\begin{equation}
  \mathrm{LP}_i = \gamma_0 + \sum_{k=1}^{P} \gamma_k z_{ik},
  \qquad
  p_i = g^{-1}(\mathrm{LP}_i),
  \label{eq:tech-risk-lp}
\end{equation}
where $\gamma_0$ is the intercept, which sets the overall level of risk, $\gamma_k$ is the weight of predictor $k$, and $g$ is commonly the logit or the complementary log-log link \cite{steyerberg2019cpm}.
The weights set the ranking of persons by risk, whereas the intercept sets the level of risk in the population the model was fitted to.

A risk model is either developed in the data at hand, with every weight estimated there, or taken from a published score whose weights were estimated elsewhere.
SCORE2 is such a score: it estimates the 10-year risk of fatal and non-fatal cardiovascular disease from age, sex, smoking, systolic blood pressure and non-HDL cholesterol, with separate calibrations for four European risk regions \cite{score2_2021}.
A published score costs no parameters in new data and carries weights estimated in a much larger population, but those weights were chosen for its own outcome and population.

\subsection{Recalibration of an Existing Score}
\label{sec:tech-recalibration}

While a published score ranks persons with weights estimated elsewhere, its predicted risks are miscalibrated in a population with a different baseline risk or for an outcome other than the one it was built for.
Updating methods address this by re-estimating part of the model in the new data, in steps that each estimate more parameters \cite{janssen2008updating, steyerberg2019cpm}.
Recalibration-in-the-large re-estimates only the intercept.
Logistic recalibration, introduced by Cox \cite{cox1958binary}, re-estimates an intercept and one slope on the published linear predictor,
\begin{equation}
  g(p_i) = a + \gamma\, \mathrm{LP}_i ,
  \label{eq:tech-recalibration}
\end{equation}
where $a$ shifts the level of risk and $\gamma$ scales all published weights by the same factor, so that the published ranking of persons is kept.
For a score reported as a risk $\hat p_i$, the published linear predictor is recovered on the scale of the link, for example $\mathrm{LP}_i = \log\{-\log(1 - \hat p_i)\}$ for the complementary log-log link.
Revision re-estimates selected weights, and refitting re-estimates every weight, so that the score's predictors are kept but its weighting is not \cite{janssen2008updating}.
Each step adapts the model more closely to the new data at the cost of more parameters, which the new data must be large enough to estimate \cite{steyerberg2019cpm}.

Recalibration keeps the relative weights of the published score, so a predictor whose effect on the new outcome differs from its effect in the score stays mis-weighted.
A block of new predictors added to a recalibrated score is then compared against this mis-weighting as well as against the score.
Correct reading of added value over a published score therefore depends on whether its relative weights fit the new outcome.
Otherwise, a new predictor that is correlated with a mis-weighted factor, such as age, can appear to add information only by correcting that weight.
